\chapter{Syntactic Transition--Output Categories as Residual Kleisli Images}
\label{ch:syntactic-transition-output-categories}

\section{Residual operations and standing conventions}
\label{sec:syn-to-standing-conventions}

The preceding chapter associated to every behaviour specification
\[
    k:K\to\mathbf{Beh}_0
\]
a derivative image
\[
    D:=\operatorname{Der}_0(K)
\]
and a minimal residual Mealy morphism
\[
    \mathsf{Min}(K):\mathbb A\rightsquigarrow\mathbb B .
\]
This chapter studies the input operations induced on this residual carrier.

We retain the notation and standing assumptions of
Chapter~\ref{ch:behaviour-categories-myhill-nerode}. Thus \(\mathcal E\) is
Barr-exact and locally cartesian closed, and
\[
    \mathbb A=(A_0,A_1,s_A,t_A,e_A,m_A),
    \qquad
    \mathbb B=(B_0,B_1,s_B,t_B,e_B,m_B)
\]
are internal categories in \(\mathcal E\), with the composition convention
\[
    m_A(a,b)=b\circ a
\]
for composable arrows
\[
    a:u\to v,
    \qquad
    b:v\to w.
\]
Thus a Mealy state over \(w\) processes the input \(b\circ a\) by first
processing \(b\), then processing \(a\).

The derivative image is written
\[
    D=\operatorname{Der}_0(K),
\]
with image factorization
\[
    D_K
        \xrightarrow{e_K}
    D
        \xrightarrow{j_K}
    \mathbf{Beh}_0.
\]
Here \(D_K\) is the formal derivative orbit object of
Chapter~\ref{ch:behaviour-categories-myhill-nerode}, and \(j_K\) is monic in
\[
    \mathcal E/(A_0\times B_0).
\]

Since \(D\) is derivative-closed, the canonical behaviour Mealy morphism
restricts to a Mealy morphism carried by
\[
    A_0
    \xleftarrow{p_D}
    D
    \xrightarrow{q_D}
    B_0 .
\]
The restricted transition and output maps are denoted
\[
    \delta_D:
    A_1\times_{t_A,A_0,p_D}D\to D,
\]
and
\[
    \omega_D:
    A_1\times_{t_A,A_0,p_D}D\to B_1.
\]
For an input arrow
\[
    a:u\to v
\]
and a residual state \(x\in D\) with \(p_D(x)=v\), we write
\[
    \partial_a x:=\delta_D(a,x).
\]
Thus
\[
    p_D(\partial_a x)=u,
\]
and the emitted output arrow is
\[
    \omega_D(a,x):q_D(\partial_a x)\to q_D(x).
\]

The residual transition and output are displayed by the diagram
\[
\begin{tikzcd}[column sep=large, row sep=large]
    q_D(\partial_a x)
        \arrow[r, "{\omega_D(a,x)}"]
    &
    q_D(x)
    \\
    \partial_a x
        \arrow[u, "q_D"]
    &
    x
        \arrow[l, dashed, "{a}"']
        \arrow[u, "q_D"']
\end{tikzcd}
\]
where the dashed arrow is not an arrow in \(D\), but the action of the input
arrow \(a\) on the residual state \(x\).

\begin{definition}[Residual transition--output operation]
\label{def:syn-to-residual-transition-output-operation}
    The \textbf{residual transition--output operation} induced by an input
    arrow
    \[
        a:u\to v
    \]
    is the assignment
    \[
        x\longmapsto
        \bigl(\partial_a x,\omega_D(a,x)\bigr),
        \qquad
        x\in D_v.
    \]
    It consists of a state transition
    \[
        \partial_a:D_v\to D_u
    \]
    together with an output-comparison arrow
    \[
        \omega_D(a,x):q_D(\partial_a x)\to q_D(x).
    \]
\end{definition}

Thus an input arrow does not merely determine a map
\[
    D_v\to D_u.
\]
It determines a map equipped with an output comparison. Consequently, the
appropriate endomorphism object is not an ordinary endomorphism category of the
indexed family \(D\), but a Kleisli endomorphism category for the
output-comparison monad determined by \(\mathbb B\).

The main construction of this chapter is the \textbf{residual Kleisli action}
\[
    \chi_D:
    \mathbb A^{\mathrm{op}}
    \longrightarrow
    \operatorname{KlEnd}_{\mathbb B}(D).
\]
The \textbf{syntactic transition--output category} of \(k\) is the regular
image of this action:
\[
    \operatorname{SynTO}_{\mathbb A,\mathbb B}(K)
    :=
    \operatorname{RegIm}(\chi_D).
\]

The relationship between the input category, the residual carrier, and the
Kleisli endomorphism category is summarized by
\[
\begin{tikzcd}[column sep=huge, row sep=large]
    \mathbb A^{\mathrm{op}}
        \arrow[r, "{\chi_D}"]
        \arrow[dr, dashed, "\operatorname{RegIm}"']
    &
    \operatorname{KlEnd}_{\mathbb B}(D)
    \\
    &
    \operatorname{SynTO}_{\mathbb A,\mathbb B}(K)
        \arrow[u, hook]
\end{tikzcd}
\]

\begin{remark}[State quotients versus operation quotients]
\label{rem:syn-to-state-versus-operation-early}
    The quotient studied in this chapter is not a quotient of residual states.
    The state quotient of Chapter~\ref{ch:behaviour-categories-myhill-nerode}
    identifies formal residual derivatives and produces the residual state
    object
    \[
        D=\operatorname{Der}_0(K).
    \]
    The quotient here identifies input arrows according to their induced
    transition--output operations on \(D\).
\end{remark}

\subsection{Running specializations}
\label{subsec:syn-to-standing-running-specializations}

In the one-object input-output case, let
\[
    \mathbb A=M,
    \qquad
    \mathbb B=N
\]
be monoid objects regarded as one-object internal categories with the convention
of Chapter~\ref{ch:spans-internal-categories-mealy}. Then \(D\) is a right
\(M\)-object, written
\[
    x\cdot m:=\partial_mx,
\]
and the output map is an \(N\)-valued cocycle
\[
    \omega_D:M\times D\to N.
\]
The residual transition--output operation of \(m\in M\) is
\[
    x\longmapsto
    \bigl(x\cdot m,\omega_D(m,x)\bigr).
\]

In the stateless case, let
\[
    F:\mathbb A\to\mathbb B
\]
be an internal functor regarded as a map-representable Mealy morphism. Its
residual states are the slice-restricted functorial behaviours
\[
    H_v^F=F\circ\operatorname{dom}_v.
\]
For an arrow \(a:u\to v\), the induced transition--output operation sends
\[
    H_v^F\longmapsto H_u^F
\]
and emits
\[
    F_1(a):F_0(u)\to F_0(v).
\]

\section{The output-comparison Kleisli category}
\label{sec:syn-to-output-comparison-kleisli-category}

The output category \(\mathbb B\) determines a monad on the slice
\[
    \mathcal E/B_0.
\]
This is the fibrewise version of the output-comparison monad \(B\circ-\) used
for Mealy morphisms in Chapter~\ref{ch:spans-internal-categories-mealy}.

\begin{definition}[Output-comparison monad]
\label{def:syn-to-output-comparison-monad}
    The \textbf{output-comparison monad} of \(\mathbb B\) is the monad
    \[
        T_{\mathbb B}:\mathcal E/B_0\to\mathcal E/B_0
    \]
    given by
    \[
        T_{\mathbb B}:=(t_B)_!(s_B)^*.
    \]
    Explicitly, for an object
    \[
        q:X\to B_0
    \]
    of \(\mathcal E/B_0\),
    \[
        T_{\mathbb B}X
        =
        X\times_{q,B_0,s_B}B_1
        \xrightarrow{t_B\pi_{B_1}}
        B_0 .
    \]
    A generalized element of \(T_{\mathbb B}X\) is a pair
    \[
        (x,\beta)
    \]
    where
    \[
        \beta:q(x)\to b
    \]
    is an arrow of \(\mathbb B\).

    The unit is induced by the identity map \(e_B:B_0\to B_1\):
    \[
        \eta_X(x)=(x,1_{q(x)}).
    \]
    The multiplication is induced by composition in \(\mathbb B\):
    \[
        \mu_X(x,\beta,\gamma)=(x,\gamma\circ\beta),
    \]
    where
    \[
        \beta:q(x)\to b,
        \qquad
        \gamma:b\to c.
    \]
\end{definition}

The pullback defining \(T_{\mathbb B}X\) is
\[
\begin{tikzcd}[column sep=large, row sep=large]
    T_{\mathbb B}X
        \arrow[r, "\pi_{B_1}"]
        \arrow[d, "\pi_X"']
        \arrow[dr, phantom, "\lrcorner", very near start]
    &
    B_1
        \arrow[d, "s_B"]
    \\
    X
        \arrow[r, "q"']
    &
    B_0 .
\end{tikzcd}
\]
Its structure map to \(B_0\) is
\[
    t_B\pi_{B_1}:T_{\mathbb B}X\to B_0.
\]

The multiplication is displayed by composing output comparisons:
\[
\begin{tikzcd}[column sep=large, row sep=large]
    q(x)
        \arrow[r, "\beta"]
        \arrow[rr, bend left=18, "\gamma\circ\beta"]
    &
    b
        \arrow[r, "\gamma"]
    &
    c .
\end{tikzcd}
\]

\begin{proposition}[Output-comparison monad]
\label{prop:syn-to-output-comparison-monad}
    The data of Definition~\ref{def:syn-to-output-comparison-monad} define a
    monad on
    \[
        \mathcal E/B_0.
    \]
\end{proposition}

\begin{proof}
    \leavevmode
    \begin{enumerate}
        \item \textbf{Check the unit laws.}
        The left and right unit laws for
        \[
            \eta_X:X\to T_{\mathbb B}X
        \]
        are precisely the left and right identity laws for the internal category
        \(\mathbb B\).

        \item \textbf{Check associativity.}
        The associativity law for
        \[
            \mu_X:T_{\mathbb B}T_{\mathbb B}X\to T_{\mathbb B}X
        \]
        is precisely associativity of composition in \(\mathbb B\). For arrows
        \[
            q(x)\xrightarrow{\beta}b
            \xrightarrow{\gamma}c
            \xrightarrow{\xi}d,
        \]
        both monad associativity composites give
        \[
            (x,\xi\circ\gamma\circ\beta).
        \]

        \item \textbf{Work in the slice.}
        All maps are maps over \(B_0\), because the structure map of
        \(T_{\mathbb B}X\) records the target of the output-comparison arrow.
        Hence the monad laws hold in \(\mathcal E/B_0\).
    \end{enumerate}
\end{proof}

\begin{definition}[Output Kleisli category]
\label{def:syn-to-output-kleisli-category}
    The \textbf{output Kleisli category} is the Kleisli category
    \[
        \operatorname{Kl}_{\mathbb B}:=\operatorname{Kl}(T_{\mathbb B}).
    \]
\end{definition}

\begin{lemma}[Kleisli normal form]
\label{lem:syn-to-kleisli-normal-form}
    Let
    \[
        q_X:X\to B_0,
        \qquad
        q_Y:Y\to B_0
    \]
    be objects of \(\mathcal E/B_0\). A Kleisli map
    \[
        X\rightsquigarrow Y
    \]
    is equivalently a pair
    \[
        F:X\to Y,
        \qquad
        \alpha:X\to B_1
    \]
    satisfying
    \[
        s_B\alpha=q_YF,
        \qquad
        t_B\alpha=q_X.
    \]
    In generalized-element notation,
    \[
        \alpha_x:q_Y(Fx)\to q_X(x).
    \]

    If
    \[
        (F,\alpha):X\rightsquigarrow Y,
        \qquad
        (G,\gamma):Y\rightsquigarrow Z,
    \]
    then their Kleisli composite is
    \[
        (G,\gamma)\odot(F,\alpha)
        =
        (GF,\alpha\star\gamma),
    \]
    where
    \[
        (\alpha\star\gamma)_x
        =
        \alpha_x\circ\gamma_{Fx}.
    \]
\end{lemma}

\begin{proof}
    \leavevmode
    \begin{enumerate}
        \item \textbf{Unpack a Kleisli map.}
        A Kleisli map
        \[
            X\rightsquigarrow Y
        \]
        is a morphism
        \[
            X\to T_{\mathbb B}Y
            =
            Y\times_{q_Y,B_0,s_B}B_1
        \]
        in the slice \(\mathcal E/B_0\). Hence it is equivalently a pair
        \[
            F:X\to Y,
            \qquad
            \alpha:X\to B_1
        \]
        satisfying
        \[
            s_B\alpha=q_YF.
        \]

        \item \textbf{Impose the slice condition.}
        The map
        \[
            X\to T_{\mathbb B}Y
        \]
        must lie over \(B_0\). Since the structure map of \(T_{\mathbb B}Y\) is
        \(t_B\pi_{B_1}\), this condition is exactly
        \[
            t_B\alpha=q_X.
        \]

        \item \textbf{Compute Kleisli composition.}
        Applying the second Kleisli map gives
        \[
            x\longmapsto
            \bigl(GF x,\gamma_{Fx},\alpha_x\bigr),
        \]
        and the multiplication of \(T_{\mathbb B}\) composes the output arrows:
        \[
        \begin{tikzcd}[column sep=large, row sep=large]
            q_Z(GFx)
                \arrow[r, "{\gamma_{Fx}}"]
                \arrow[rr, bend left=18, "{\alpha_x\circ\gamma_{Fx}}"]
            &
            q_Y(Fx)
                \arrow[r, "{\alpha_x}"]
            &
            q_X(x).
        \end{tikzcd}
        \]
        Thus the composite output component is
        \[
            (\alpha\star\gamma)_x
            =
            \alpha_x\circ\gamma_{Fx}.
        \]
    \end{enumerate}
\end{proof}

\begin{remark}[Composition order]
\label{rem:syn-to-kleisli-composition-order}
    With the present output-comparison orientation, the output attached to the
    first Kleisli map is composed after the output attached to the second state
    map:
    \[
        q_Z(GFx)
            \xrightarrow{\gamma_{Fx}}
        q_Y(Fx)
            \xrightarrow{\alpha_x}
        q_X(x).
    \]
    Hence the composite output is
    \[
        \alpha_x\circ\gamma_{Fx},
    \]
    not \(\gamma_{Fx}\circ\alpha_x\).
\end{remark}

\subsection{Running specializations}
\label{subsec:syn-to-output-kleisli-running-specializations}

If \(\mathbb B=N\) is the one-object internal category associated to a monoid
object \(N\), then
\[
    B_0=1,
    \qquad
    B_1=N,
\]
and
\[
    T_{\mathbb B}X=X\times N.
\]
Thus \(\operatorname{Kl}_{\mathbb B}\) is the usual writer Kleisli category.
A Kleisli map
\[
    X\rightsquigarrow Y
\]
is a pair
\[
    F:X\to Y,
    \qquad
    o:X\to N.
\]
Composition is
\[
    (G,o')\odot(F,o)
    =
    \bigl(GF,\;x\mapsto o(x)o'(Fx)\bigr),
\]
using the one-object convention from Chapter~\ref{ch:spans-internal-categories-mealy}.

If \(\mathbb B\) is chaotic on an object \(B_0\), then an output comparison is
uniquely determined by its source and target. In that case a Kleisli map is
essentially a state map together with its endpoint typing; there is no
additional independent output data.

\section{Kleisli endomorphisms of a residual family}
\label{sec:syn-to-kleisli-endomorphisms-residual-family}

The carrier of the minimal residual realization is a map
\[
    D\to A_0\times B_0.
\]
Equivalently, it is an \(A_0\)-indexed family of objects of
\(\mathcal E/B_0\):
\[
    v\longmapsto D_v\to B_0.
\]
For each \(v\in A_0\), the fibre \(D_v\) is therefore an object of the output
Kleisli category
\[
    \operatorname{Kl}_{\mathbb B}.
\]

\begin{definition}[Kleisli endomorphism category]
\label{def:syn-to-kleisli-endomorphism-category}
    The \textbf{Kleisli endomorphism category} of the residual family \(D\) is
    the internal category
    \[
        \operatorname{KlEnd}_{\mathbb B}(D)
    \]
    whose object object is \(A_0\), and whose arrows
    \[
        v\to u
    \]
    are Kleisli maps
    \[
        D_v\rightsquigarrow D_u
    \]
    in
    \[
        \operatorname{Kl}_{\mathbb B}.
    \]
\end{definition}

Thus an arrow of \(\operatorname{KlEnd}_{\mathbb B}(D)\) from \(v\) to \(u\)
has the form
\[
\begin{tikzcd}[column sep=huge, row sep=large]
    D_v
        \arrow[r, dashed, "{(F,\alpha)}"]
        \arrow[dr, "q_D"']
    &
    D_u
        \arrow[d, "q_D"]
    \\
    &
    B_0
\end{tikzcd}
\]
where the dashed arrow denotes a Kleisli map. In normal form it is a pair
\[
    F:D_v\to D_u,
    \qquad
    \alpha_x:q_D(Fx)\to q_D(x).
\]

\begin{theorem}[Existence of the Kleisli endomorphism category]
\label{thm:syn-to-kleisli-endomorphism-category-exists}
    The fibrewise Kleisli endomorphisms of \(D\) are represented by an internal
    category
    \[
        \operatorname{KlEnd}_{\mathbb B}(D)
    \]
    with object object \(A_0\).
\end{theorem}

\begin{proof}
    \leavevmode
    \begin{enumerate}
        \item \textbf{Parametrize source and target fibres.}
        Put
        \[
            P:=A_0\times A_0,
        \]
        and write
        \[
            \sigma,\tau:P\to A_0
        \]
        for the projections
        \[
            \sigma(v,u)=v,
            \qquad
            \tau(v,u)=u.
        \]
        Pull back the residual family \(D\to A_0\times B_0\) along these
        projections:
        \[
            D_{\mathrm{src}}:=P\times_{\sigma,A_0,p_D}D,
            \qquad
            D_{\mathrm{tgt}}:=P\times_{\tau,A_0,p_D}D.
        \]
        Over a generalized point \((v,u)\in P\), these are
        \[
            (D_{\mathrm{src}})_{(v,u)}=D_v,
            \qquad
            (D_{\mathrm{tgt}})_{(v,u)}=D_u.
        \]

        The situation is displayed by
        \[
        \begin{tikzcd}[column sep=large, row sep=large]
            D_{\mathrm{src}}
                \arrow[r]
                \arrow[d]
                \arrow[dr, phantom, "\lrcorner", very near start]
            &
            D
                \arrow[d, "p_D"]
            \\
            A_0\times A_0
                \arrow[r, "\sigma"']
            &
            A_0
        \end{tikzcd}
        \qquad
        \begin{tikzcd}[column sep=large, row sep=large]
            D_{\mathrm{tgt}}
                \arrow[r]
                \arrow[d]
                \arrow[dr, phantom, "\lrcorner", very near start]
            &
            D
                \arrow[d, "p_D"]
            \\
            A_0\times A_0
                \arrow[r, "\tau"']
            &
            A_0 .
        \end{tikzcd}
        \]

        \item \textbf{Construct the fibrewise Kleisli target.}
        Both \(D_{\mathrm{src}}\) and \(D_{\mathrm{tgt}}\) carry maps to
        \(P\times B_0\), using their projection to \(P\) and the map \(q_D\).
        Form the fibrewise output-comparison object
        \[
            T_{\mathbb B}^{P}(D_{\mathrm{tgt}})
            =
            D_{\mathrm{tgt}}
            \times_{q_D,B_0,s_B}
            B_1.
        \]
        Its map to \(P\times B_0\) is given by the projection to \(P\) and by
        \(t_B\) on the \(B_1\)-component.

        The defining pullback is
        \[
        \begin{tikzcd}[column sep=large, row sep=large]
            T_{\mathbb B}^{P}(D_{\mathrm{tgt}})
                \arrow[r, "\pi_{B_1}"]
                \arrow[d, "\pi_{D_{\mathrm{tgt}}}"']
                \arrow[dr, phantom, "\lrcorner", very near start]
            &
            B_1
                \arrow[d, "s_B"]
            \\
            D_{\mathrm{tgt}}
                \arrow[r, "q_D"']
            &
            B_0 .
        \end{tikzcd}
        \]

        \item \textbf{Represent Kleisli maps by dependent products.}
        A Kleisli map over \(P\)
        \[
            D_{\mathrm{src}}\rightsquigarrow D_{\mathrm{tgt}}
        \]
        is a map over \(P\times B_0\)
        \[
            D_{\mathrm{src}}\to T_{\mathbb B}^{P}(D_{\mathrm{tgt}}).
        \]
        Equivalently, it is a section of the pullback of
        \(T_{\mathbb B}^{P}(D_{\mathrm{tgt}})\) along
        \[
            D_{\mathrm{src}}\to P\times B_0.
        \]
        More explicitly, the relevant object over \(D_{\mathrm{src}}\) is
        \[
            D_{\mathrm{src}}
            \times_{P\times B_0}
            T_{\mathbb B}^{P}(D_{\mathrm{tgt}}),
        \]
        where \(D_{\mathrm{src}}\to P\times B_0\) uses its projection to \(P\)
        and \(q_D\), and
        \(T_{\mathbb B}^{P}(D_{\mathrm{tgt}})\to P\times B_0\) uses its
        projection to \(P\) and the target of the output-comparison arrow.
        Since \(\mathcal E\) is locally cartesian closed, the dependent product
        along
        \[
            D_{\mathrm{src}}\to P
        \]
        represents these sections:
        \[
            \Pi_{D_{\mathrm{src}}\to P}
            \left(
                D_{\mathrm{src}}
                \times_{P\times B_0}
                T_{\mathbb B}^{P}(D_{\mathrm{tgt}})
            \right)
            \to P.
        \]
        This gives an object
        \[
            \operatorname{KlEnd}_{\mathbb B}(D)_1\to P.
        \]

        \item \textbf{Define source, target, identity, and composition.}
        The source and target maps
        \[
            \operatorname{KlEnd}_{\mathbb B}(D)_1\rightrightarrows A_0
        \]
        are induced by the two projections
        \[
            P=A_0\times A_0\rightrightarrows A_0.
        \]
        Identities are represented by the Kleisli units
        \[
            D_v\to T_{\mathbb B}D_v.
        \]
        Composition is represented by fibrewise Kleisli composition.

        \item \textbf{Verify the internal-category laws.}
        The unit and associativity laws follow from the unit and associativity
        laws of the Kleisli category
        \[
            \operatorname{Kl}_{\mathbb B}.
        \]
        Therefore the represented data form an internal category.
    \end{enumerate}
\end{proof}

\begin{theorem}[Transition--output normal form]
\label{thm:syn-to-transition-output-normal-form}
    An arrow
    \[
        v\to u
    \]
    of
    \[
        \operatorname{KlEnd}_{\mathbb B}(D)
    \]
    is equivalently a pair
    \[
        (F,\alpha)
    \]
    consisting of a map
    \[
        F:D_v\to D_u
    \]
    and, for each \(x\in D_v\), an output-comparison arrow
    \[
        \alpha_x:q_D(Fx)\to q_D(x)
    \]
    in \(\mathbb B\).

    If
    \[
        (F,\alpha):v\to u,
        \qquad
        (G,\gamma):u\to w,
    \]
    then their composite is
    \[
        (G,\gamma)\circ(F,\alpha)
        =
        (GF,\alpha\star\gamma),
    \]
    where
    \[
        (\alpha\star\gamma)_x
        =
        \alpha_x\circ\gamma_{Fx}.
    \]
\end{theorem}

\begin{proof}
    \leavevmode
    \begin{enumerate}
        \item \textbf{Apply Kleisli normal form fibrewise.}
        This is Lemma~\ref{lem:syn-to-kleisli-normal-form} applied to the
        objects
        \[
            D_v\to B_0,
            \qquad
            D_u\to B_0
        \]
        of \(\mathcal E/B_0\).

        \item \textbf{Compute composition.}
        The formula for composition is the fibrewise Kleisli composition
        formula. The output-comparison part is the composite
        \[
        \begin{tikzcd}[column sep=large, row sep=large]
            q_D(GFx)
                \arrow[r, "{\gamma_{Fx}}"]
                \arrow[rr, bend left=18, "{\alpha_x\circ\gamma_{Fx}}"]
            &
            q_D(Fx)
                \arrow[r, "{\alpha_x}"]
            &
            q_D(x).
        \end{tikzcd}
        \]
    \end{enumerate}
\end{proof}

\subsection{Running specializations}
\label{subsec:syn-to-kleisli-endomorphism-running-specializations}

If \(\mathbb B=N\) is one-object but \(A_0\) is arbitrary, then an arrow
\[
    v\to u
\]
of \(\operatorname{KlEnd}_{N}(D)\) is a pair
\[
    F:D_v\to D_u,
    \qquad
    o:D_v\to N.
\]
Thus, fibrewise,
\[
    \operatorname{KlEnd}_{N}(D)(v,u)
    \cong
    D_u^{D_v}\times N^{D_v}.
\]

If both \(\mathbb A=M\) and \(\mathbb B=N\) are one-object, then
\(A_0=1\), the residual family is an unindexed object \(D\), and
\[
    \operatorname{KlEnd}_{N}(D)
\]
is the one-object internal category whose arrow monoid is
\[
    D^D\ltimes N^D.
\]
An arrow is a pair
\[
    (\tau,o),
    \qquad
    \tau:D\to D,
    \qquad
    o:D\to N,
\]
and composition is
\[
    (\tau,o)(\tau',o')
    =
    \left(
        \tau'\circ\tau,\,
        x\longmapsto o(x)o'(\tau x)
    \right).
\]
This is the writer--wreath form of the Kleisli endomorphism monoid.

\section{The residual Kleisli action}
\label{sec:syn-to-residual-kleisli-action}

The Mealy structure on \(D\) is a Kleisli action of
\(\mathbb A^{\mathrm{op}}\) on the indexed family \(D\). The passage to
\(\mathbb A^{\mathrm{op}}\) accounts for the fact that an input arrow
\[
    a:u\to v
\]
acts contravariantly on residual states:
\[
    D_v\rightsquigarrow D_u.
\]

Over the arrow object \(A_1\), the two relevant pulled-back families are
\[
    t_A^*D
    \qquad\text{and}\qquad
    s_A^*D.
\]
The transition-output structure is a Kleisli morphism
\[
    t_A^*D\rightsquigarrow s_A^*D
\]
over \(A_1\), represented by
\[
    (a,x)
    \longmapsto
    \bigl(\partial_a x,\omega_D(a,x)\bigr).
\]

The fibrewise action is displayed by
\[
\begin{tikzcd}[column sep=large, row sep=large]
    D_v
        \arrow[r, dashed, "{\chi_D(a)}"]
        \arrow[dr, "q_D"']
    &
    D_u
        \arrow[d, "q_D"]
    \\
    &
    B_0
\end{tikzcd}
\qquad
\begin{tikzcd}[column sep=large, row sep=large]
    q_D(\partial_a x)
        \arrow[r, "{\omega_D(a,x)}"]
    &
    q_D(x).
\end{tikzcd}
\]

\begin{definition}[Residual Kleisli action]
\label{def:syn-to-residual-kleisli-action}
    The \textbf{residual Kleisli action} of
    \[
        D=\operatorname{Der}_0(K)
    \]
    is the internal functor
    \[
        \chi_D:
        \mathbb A^{\mathrm{op}}
        \longrightarrow
        \operatorname{KlEnd}_{\mathbb B}(D)
    \]
    defined as follows.

    On objects, \(\chi_D\) is the identity map
    \[
        A_0\to A_0.
    \]

    For an arrow
    \[
        a:u\to v
    \]
    of \(\mathbb A\), regarded as an arrow
    \[
        a^{\mathrm{op}}:v\to u
    \]
    of \(\mathbb A^{\mathrm{op}}\), the arrow
    \[
        \chi_D(a):v\to u
    \]
    is the Kleisli map
    \[
        D_v\rightsquigarrow D_u
    \]
    represented by
    \[
        x\longmapsto
        \bigl(\partial_a x,\omega_D(a,x)\bigr).
    \]
    In transition--output normal form,
    \[
        \chi_D(a)=(F_a,\alpha_a),
    \]
    where
    \[
        F_a(x)=\partial_a x,
        \qquad
        (\alpha_a)_x=\omega_D(a,x).
    \]
\end{definition}

\begin{proposition}[Mealy laws as Kleisli action laws]
\label{prop:syn-to-mealy-laws-kleisli-action-laws}
    The Mealy unit and multiplication laws for the minimal residual Mealy
    morphism \(\mathsf{Min}(K)\) are precisely the unit and associativity laws
    for the Kleisli action of \(\mathbb A^{\mathrm{op}}\) on \(D\).
\end{proposition}

\begin{proof}
    \leavevmode
    \begin{enumerate}
        \item \textbf{Check the unit law.}
        The identity input \(1_v\) acts by
        \[
            x\longmapsto
            \bigl(\partial_{1_v}x,\omega_D(1_v,x)\bigr).
        \]
        By the Mealy unit laws this is
        \[
            x\longmapsto
            (x,1_{q_Dx}),
        \]
        which is the identity Kleisli endomorphism of \(D_v\).

        \item \textbf{Check the multiplication law.}
        For composable arrows
        \[
            a:u\to v,
            \qquad
            b:v\to w,
        \]
        and \(x\in D_w\), the Kleisli composite of the actions of \(b\) and
        \(a\) sends \(x\) to
        \[
            \left(
                \partial_a(\partial_bx),
                \omega_D(b,x)\circ\omega_D(a,\partial_bx)
            \right).
        \]
        The Mealy multiplication laws say exactly that this is
        \[
            \left(
                \partial_{b\circ a}x,
                \omega_D(b\circ a,x)
            \right),
        \]
        which is the action of the composite input \(b\circ a\).
    \end{enumerate}
\end{proof}

\begin{theorem}[Functoriality of the residual Kleisli action]
\label{thm:syn-to-residual-kleisli-action-functorial}
    The assignment of Definition~\ref{def:syn-to-residual-kleisli-action}
    defines an internal functor
    \[
        \chi_D:
        \mathbb A^{\mathrm{op}}
        \to
        \operatorname{KlEnd}_{\mathbb B}(D).
    \]
\end{theorem}

\begin{proof}
    \leavevmode
    \begin{enumerate}
        \item \textbf{Preserve identities.}
        For \(x\in D_v\),
        \[
            \partial_{1_v}x=x
        \]
        and
        \[
            \omega_D(1_v,x)=1_{q_Dx}.
        \]
        Therefore \(\chi_D(1_v)\) is the identity Kleisli endomorphism of
        \(D_v\).

        \item \textbf{Preserve composition.}
        Let
        \[
            a:u\to v,
            \qquad
            b:v\to w
        \]
        be composable arrows in \(\mathbb A\), and let \(x\in D_w\). In
        \(\mathbb A^{\mathrm{op}}\), the composite of
        \[
            b^{\mathrm{op}}:w\to v
            \qquad\text{and}\qquad
            a^{\mathrm{op}}:v\to u
        \]
        is
        \[
            (b\circ a)^{\mathrm{op}}:w\to u.
        \]
        Hence we must prove
        \[
            \chi_D(b\circ a)
            =
            \chi_D(a)\circ\chi_D(b).
        \]

        The composite Kleisli operation \(\chi_D(a)\circ\chi_D(b)\) sends
        \(x\) to
        \[
            \left(
                \partial_a(\partial_bx),
                \omega_D(b,x)\circ\omega_D(a,\partial_bx)
            \right).
        \]
        By the Mealy multiplication laws this is
        \[
            \left(
                \partial_{b\circ a}x,
                \omega_D(b\circ a,x)
            \right),
        \]
        which is exactly \(\chi_D(b\circ a)(x)\).

        The composite action is displayed by
        \[
        \begin{tikzcd}[column sep=large, row sep=large]
            D_w
                \arrow[r, dashed, "{\chi_D(b)}"]
                \arrow[rr, dashed, bend left=16, "{\chi_D(b\circ a)}"]
            &
            D_v
                \arrow[r, dashed, "{\chi_D(a)}"]
            &
            D_u .
        \end{tikzcd}
        \]
    \end{enumerate}
\end{proof}

\subsection{Running specializations}
\label{subsec:syn-to-residual-action-running-specializations}

In the one-object input-output case, the residual Kleisli action is a monoid
morphism
\[
    \chi_D:
    M\to D^D\ltimes N^D.
\]
It sends \(m\in M\) to
\[
    \chi_D(m)=(\tau_m,o_m),
\]
where
\[
    \tau_m(x)=x\cdot m,
    \qquad
    o_m(x)=\omega_D(m,x).
\]
Strictly, the one-object category \(\mathbb A^{\mathrm{op}}\) corresponds to
the opposite categorical composition. Throughout this specialization we use the
right-action identification of \(\mathbb A^{\mathrm{op}}\) with the underlying
input monoid \(M\), so that the categorical functoriality law for \(\chi_D\) is
written as the usual right-action law
\[
    x\cdot(mn)=(x\cdot m)\cdot n.
\]
The homomorphism law is exactly the right action law together with the output
cocycle law:
\[
    \omega_D(mn,x)
    =
    \omega_D(m,x)\omega_D(n,x\cdot m).
\]

In the stateless case, for an internal functor
\[
    F:\mathbb A\to\mathbb B,
\]
the residual action sends
\[
    H_v^F
        \longmapsto
    H_u^F
\]
along an arrow
\[
    a:u\to v,
\]
and its output component is
\[
    F_1(a):F_0(u)\to F_0(v).
\]
Thus the action of \(\mathbb A^{\mathrm{op}}\) records both slice
reindexing and the arrow map of \(F\).

\section{The syntactic transition--output image}
\label{sec:syn-to-syntactic-transition-output-image}

We now define the central object of the chapter.

\begin{definition}[Syntactic transition--output category]
\label{def:syn-to-syntactic-transition-output-category}
    Let
    \[
        k:K\to\mathbf{Beh}_0
    \]
    be a behaviour specification, and put
    \[
        D:=\operatorname{Der}_0(K).
    \]
    The \textbf{syntactic transition--output category} of \(k\) is the regular
    image of the residual Kleisli action
    \[
        \chi_D:
        \mathbb A^{\mathrm{op}}
        \to
        \operatorname{KlEnd}_{\mathbb B}(D).
    \]
    It is denoted
    \[
        \operatorname{SynTO}_{\mathbb A,\mathbb B}(K).
    \]
\end{definition}

Since \(\chi_D\) is identity on objects, the object object of
\[
    \operatorname{SynTO}_{\mathbb A,\mathbb B}(K)
\]
is \(A_0\). Its arrow object is the regular image of the arrow map
\[
    A_1\to\operatorname{KlEnd}_{\mathbb B}(D)_1
\]
in the slice
\[
    \mathcal E/(A_0\times A_0),
\]
where \(A_1\) is regarded as an object over \(A_0\times A_0\) by
\[
    (t_A,s_A):A_1\to A_0\times A_0.
\]
This is the correct source--target map because an arrow
\[
    a:u\to v
\]
of \(\mathbb A\) is an arrow
\[
    a^{\mathrm{op}}:v\to u
\]
of \(\mathbb A^{\mathrm{op}}\).

Thus the image factorization has the form
\[
    A_1
        \xrightarrow{e_{\mathrm{syn}}}
    \operatorname{SynTO}_{\mathbb A,\mathbb B}(K)_1
        \xrightarrow{m_{\mathrm{syn}}}
    \operatorname{KlEnd}_{\mathbb B}(D)_1
\]
in
\[
    \mathcal E/(A_0\times A_0),
\]
where \(e_{\mathrm{syn}}\) is a regular epimorphism and
\(m_{\mathrm{syn}}\) is monic.

The image factorization is displayed by
\[
\begin{tikzcd}[column sep=large, row sep=large]
    A_1
        \arrow[rr, "{\chi_D}"]
        \arrow[dr, two heads, "{e_{\mathrm{syn}}}"']
    &&
    \operatorname{KlEnd}_{\mathbb B}(D)_1
    \\
    &
    \operatorname{SynTO}_{\mathbb A,\mathbb B}(K)_1
        \arrow[ur, hook, "{m_{\mathrm{syn}}}"']
    &
\end{tikzcd}
\]
in the slice \(\mathcal E/(A_0\times A_0)\).

\begin{theorem}[Image subcategory]
\label{thm:syn-to-image-subcategory}
    The regular image of
    \[
        \chi_D:
        \mathbb A^{\mathrm{op}}
        \to
        \operatorname{KlEnd}_{\mathbb B}(D)
    \]
    is an internal subcategory of
    \[
        \operatorname{KlEnd}_{\mathbb B}(D).
    \]
    Hence \(\chi_D\) factors as an internal functor
    \[
        \mathbb A^{\mathrm{op}}
            \twoheadrightarrow
        \operatorname{SynTO}_{\mathbb A,\mathbb B}(K)
            \hookrightarrow
        \operatorname{KlEnd}_{\mathbb B}(D),
    \]
    where the first functor is identity on objects and regular epimorphic on
    arrows.
\end{theorem}

\begin{proof}
    \leavevmode
    \begin{enumerate}
        \item \textbf{Take the arrow image.}
        The object part of \(\chi_D\) is the identity on \(A_0\). The arrow part
        has a regular epi--mono factorization in the exact slice
        \[
            \mathcal E/(A_0\times A_0).
        \]
        This gives the arrow object of the proposed image subcategory.

        \item \textbf{Check identities.}
        For each \(v\in A_0\), the identity arrow of
        \(\operatorname{KlEnd}_{\mathbb B}(D)\) at \(v\) is
        \[
            \chi_D(1_v).
        \]
        Hence the identity map factors through the image subobject.

        This factorization is represented by
\[
\begin{tikzcd}[column sep=large, row sep=large]
    A_0
        \arrow[r, "e_A"]
        \arrow[dr, dashed, "\eta_{\operatorname{SynTO}}"']
    &
    A_1
        \arrow[d, two heads, "e_{\mathrm{syn}}"]
        \arrow[r, "{\chi_D}"]
    &
    \operatorname{KlEnd}_{\mathbb B}(D)_1
    \\
    &
    \operatorname{SynTO}_{\mathbb A,\mathbb B}(K)_1
        \arrow[ur, hook, "{m_{\mathrm{syn}}}"']
    &
\end{tikzcd}
\]

        \item \textbf{Check closure under composition.}
        Let
        \[
            \operatorname{SynTO}_2
            =
            \operatorname{SynTO}_1
            \times_{A_0}
            \operatorname{SynTO}_1
        \]
        be the object of composable pairs in the image. Since
        \[
            e_{\mathrm{syn}}:A_1\twoheadrightarrow\operatorname{SynTO}_1
        \]
        is a regular epimorphism in
        \[
            \mathcal E/(A_0\times A_0),
        \]
        the induced map on composable-pair objects
        \[
            A^{\mathrm{op}}_2
            =
            A_1\times_{A_0}A_1
            \longrightarrow
            \operatorname{SynTO}_1\times_{A_0}\operatorname{SynTO}_1
            =
            \operatorname{SynTO}_2
        \]
        is a regular epimorphism. This follows from stability of regular
        epimorphisms under pullback and composition.

        The situation is
        \[
        \begin{tikzcd}[column sep=large, row sep=large]
            A^{\mathrm{op}}_2
                \arrow[r, two heads]
                \arrow[d, "m_{A^{\mathrm{op}}}"']
            &
            \operatorname{SynTO}_2
                \arrow[d, dashed, "m_{\operatorname{SynTO}}"]
            \\
            A_1
                \arrow[r, two heads, "e_{\mathrm{syn}}"']
            &
            \operatorname{SynTO}_1 .
        \end{tikzcd}
        \]

        After this pullback, a composable pair in the image is represented by a
        composable pair in \(\mathbb A^{\mathrm{op}}\). Since \(\chi_D\) is an
        internal functor, the composite in
        \(\operatorname{KlEnd}_{\mathbb B}(D)\) is the image under \(\chi_D\)
        of the composite in \(\mathbb A^{\mathrm{op}}\). Therefore, after
        precomposition with a regular epimorphism, the composite lands in the
        image subobject.

        \item \textbf{Descend through the regular epimorphism.}
        By regular-epi descent for subobjects, the composite itself factors
        through the image subobject. Thus the image is closed under identities
        and composition, and hence is an internal subcategory.
    \end{enumerate}
\end{proof}

\begin{definition}[Syntactic transition--output congruence]
\label{def:syn-to-syntactic-congruence}
    The \textbf{syntactic transition--output congruence} of \(k\) is the kernel
    pair
    \[
        {\equiv_K^{\mathrm{syn}}}
        \rightrightarrows
        A_1
    \]
    of the arrow map
    \[
        \chi_D:
        A_1\to
        \operatorname{KlEnd}_{\mathbb B}(D)_1
    \]
    in the slice
    \[
        \mathcal E/(A_0\times A_0).
    \]
\end{definition}

The kernel-pair square is
\[
\begin{tikzcd}[column sep=large, row sep=large]
    {\equiv_K^{\mathrm{syn}}}
        \arrow[r, "\pi_2"]
        \arrow[d, "\pi_1"']
        \arrow[dr, phantom, "\lrcorner", very near start]
    &
    A_1
        \arrow[d, "\chi_D"]
    \\
    A_1
        \arrow[r, "\chi_D"']
    &
    \operatorname{KlEnd}_{\mathbb B}(D)_1 .
\end{tikzcd}
\]
Since the kernel pair is taken in
\[
    \mathcal E/(A_0\times A_0),
\]
it compares only parallel arrows of \(\mathbb A\).

\begin{theorem}[Syntactic quotient theorem]
\label{thm:syn-to-syntactic-quotient}
    The quotient internal category
    \[
        \mathbb A^{\mathrm{op}}/{\equiv_K^{\mathrm{syn}}}
    \]
    exists and is canonically isomorphic to
    \[
        \operatorname{SynTO}_{\mathbb A,\mathbb B}(K).
    \]
\end{theorem}

\begin{proof}
    \leavevmode
    \begin{enumerate}
        \item \textbf{Use exactness of the slice.}
        Since \(\mathcal E\) is Barr-exact, the slice
        \[
            \mathcal E/(A_0\times A_0)
        \]
        is Barr-exact. Hence the quotient of the kernel pair of
        \[
            \chi_D:
            A_1\to\operatorname{KlEnd}_{\mathbb B}(D)_1
        \]
        exists.

        \item \textbf{Identify the quotient with the regular image.}
        In a Barr-exact category, the quotient of the kernel pair of a morphism
        is its regular image. Therefore the quotient arrow object is precisely
        \[
            \operatorname{SynTO}_{\mathbb A,\mathbb B}(K)_1.
        \]

        \item \textbf{Descend identities and composition.}
        The identity and composition maps descend because
        \[
            \chi_D
        \]
        is an internal functor. The descended internal category structure is
        exactly the image subcategory structure of
        Theorem~\ref{thm:syn-to-image-subcategory}.
    \end{enumerate}
\end{proof}

\begin{proposition}[Normal form of the syntactic congruence]
\label{prop:syn-to-congruence-normal-form}
    Let
    \[
        a,a':u\to v
    \]
    be parallel arrows of \(\mathbb A\). Then
    \[
        a\equiv_K^{\mathrm{syn}}a'
    \]
    if and only if the induced Kleisli maps
    \[
        D_v\rightsquigarrow D_u
    \]
    are equal. Equivalently, for every \(x\in D_v\),
    \[
        \partial_a x=\partial_{a'}x,
    \]
    and
    \[
        \omega_D(a,x)=\omega_D(a',x).
    \]
\end{proposition}

\begin{proof}
    \leavevmode
    \begin{enumerate}
        \item \textbf{Use the definition of the kernel pair.}
        By definition,
        \[
            a\equiv_K^{\mathrm{syn}}a'
        \]
        precisely when
        \[
            \chi_D(a)=\chi_D(a')
        \]
        as arrows of
        \[
            \operatorname{KlEnd}_{\mathbb B}(D).
        \]

        \item \textbf{Apply transition--output normal form.}
        Equality of Kleisli maps in transition--output normal form is exactly
        equality of both the transition part and the output-comparison part:
        \[
            \partial_a x=\partial_{a'}x,
            \qquad
            \omega_D(a,x)=\omega_D(a',x)
        \]
        for every \(x\in D_v\). The equality of output arrows is well-typed
        because equality of the transition parts gives
        \[
            q_D(\partial_a x)=q_D(\partial_{a'}x),
        \]
        so the two output-comparison arrows have the same source and target.
    \end{enumerate}
\end{proof}

\subsection{Running specializations}
\label{subsec:syn-to-image-running-specializations}

In the one-object input-output case, the syntactic transition--output category
is the one-object category whose arrow monoid is
\[
    \operatorname{im}
    \left(
        M\to D^D\ltimes N^D
    \right).
\]
The syntactic congruence is
\[
    m\equiv_K^{\mathrm{syn}}n
\]
if and only if, for every \(x\in D\),
\[
    x\cdot m=x\cdot n,
\]
and
\[
    \omega_D(m,x)=\omega_D(n,x).
\]
Thus it remembers both the residual state transformation and the cumulative
output emitted from every residual state.

In the stateless case, the syntactic transition--output category is the image
of
\[
    \mathbb A^{\mathrm{op}}
\]
acting on the residual restrictions of an internal functor
\[
    F:\mathbb A\to\mathbb B.
\]
Two parallel input arrows are identified exactly when they induce the same
reindexing operation on the generated residual functors and the same emitted
arrow under \(F\).

\section{Universal property and context expansion}
\label{sec:syn-to-universal-property-context-expansion}

The syntactic transition--output category is the coarsest regular quotient of
the input category through which the residual Kleisli action factors. It also
admits a contextual description analogous to the classical two-sided syntactic
congruence.

\begin{definition}[Transition--output quotient]
\label{def:syn-to-transition-output-quotient}
    Let
    \[
        D=\operatorname{Der}_0(K).
    \]
    A \textbf{transition--output quotient realizing the residual action} is a
    factorization
    \[
        \mathbb A^{\mathrm{op}}
            \xrightarrow{q}
        \mathbb C
            \xrightarrow{j}
        \operatorname{KlEnd}_{\mathbb B}(D)
    \]
    such that
    \[
        q_0=1_{A_0},
    \]
    the arrow map
    \[
        q_1:A_1\twoheadrightarrow C_1
    \]
    is a regular epimorphism in
    \[
        \mathcal E/(A_0\times A_0),
    \]
    and
    \[
        \chi_D=jq.
    \]
\end{definition}

The factorization is displayed by
\[
\begin{tikzcd}[column sep=huge, row sep=large]
    \mathbb A^{\mathrm{op}}
        \arrow[r, "q", two heads]
        \arrow[dr, "{\chi_D}"']
    &
    \mathbb C
        \arrow[d, "j"]
    \\
    &
    \operatorname{KlEnd}_{\mathbb B}(D).
\end{tikzcd}
\]

\begin{theorem}[Universal property of the syntactic image]
\label{thm:syn-to-universal-property}
    Let
    \[
        \mathbb A^{\mathrm{op}}
            \xrightarrow{q}
        \mathbb C
            \xrightarrow{j}
        \operatorname{KlEnd}_{\mathbb B}(D)
    \]
    be a transition--output quotient realizing the residual action on
    \(D=\operatorname{Der}_0(K)\). Then there is a unique internal functor
    \[
        \bar q:
        \mathbb C\to
        \operatorname{SynTO}_{\mathbb A,\mathbb B}(K)
    \]
    identity on objects such that
    \[
        \bar q q=e_{\mathrm{syn}},
    \]
    where
    \[
        e_{\mathrm{syn}}:
        \mathbb A^{\mathrm{op}}
        \twoheadrightarrow
        \operatorname{SynTO}_{\mathbb A,\mathbb B}(K)
    \]
    is the image quotient. Moreover, the arrow map
    \[
        \bar q_1:
        C_1\to
        \operatorname{SynTO}_{\mathbb A,\mathbb B}(K)_1
    \]
    is a regular epimorphism.
\end{theorem}

\begin{proof}
    \leavevmode
    \begin{enumerate}
        \item \textbf{Work in the exact slice.}
        Work in
        \[
            \mathcal E/(A_0\times A_0).
        \]
        Since
        \[
            \chi_D=jq,
        \]
        every pair of arrows identified by \(q_1\) is also identified by
        \(\chi_D\). Hence
        \[
            \ker(q_1)\subseteq\ker(\chi_D).
        \]

        \item \textbf{Descend the image quotient.}
        The image quotient
        \[
            e_{\mathrm{syn},1}:A_1\twoheadrightarrow
            \operatorname{SynTO}_{\mathbb A,\mathbb B}(K)_1
        \]
        is the quotient of the kernel pair of \(\chi_D\). Therefore it
        coequalizes the kernel pair of \(q_1\). Since \(q_1\) is the effective
        quotient of its kernel pair, \(e_{\mathrm{syn},1}\) descends uniquely
        through \(q_1\), giving a unique map
        \[
            \bar q_1:
            C_1\to
            \operatorname{SynTO}_{\mathbb A,\mathbb B}(K)_1
        \]
        satisfying
        \[
            \bar q_1q_1=e_{\mathrm{syn},1}.
        \]

        \item \textbf{Verify functoriality.}
        The object map is forced to be \(1_{A_0}\). Preservation of identities
        holds after precomposition with \(q_0=1_{A_0}\) and \(q_1\), because
        both \(q\) and \(e_{\mathrm{syn}}\) are internal functors. Preservation
        of composition holds after precomposition with the induced regular
        epimorphism
        \[
            q_2:A^{\mathrm{op}}_2\twoheadrightarrow C_2.
        \]
        Since \(q_2\) is epic, the descended map preserves composition.
        Therefore \(\bar q\) is an internal functor.

        \item \textbf{Show regular epimorphy.}
        Since
        \[
            e_{\mathrm{syn},1}=\bar q_1q_1
        \]
        is a regular epimorphism and \(q_1\) is a regular epimorphism, the
        regular image of \(\bar q_1\) is all of
        \[
            \operatorname{SynTO}_{\mathbb A,\mathbb B}(K)_1.
        \]
        Hence \(\bar q_1\) is a regular epimorphism.

        \item \textbf{Display the universal triangle.}
        The construction is summarized by
        \[
        \begin{tikzcd}[column sep=huge, row sep=large]
            \mathbb A^{\mathrm{op}}
                \arrow[r, two heads, "q"]
                \arrow[d, two heads, "e_{\mathrm{syn}}"']
            &
            \mathbb C
                \arrow[dl, dashed, "\bar q"]
                \arrow[d, "j"]
            \\
            \operatorname{SynTO}_{\mathbb A,\mathbb B}(K)
                \arrow[r, hook]
            &
            \operatorname{KlEnd}_{\mathbb B}(D).
        \end{tikzcd}
        \]
    \end{enumerate}
\end{proof}

\begin{corollary}[Smallest regular quotient]
\label{cor:syn-to-smallest-regular-quotient}
    The internal category
    \[
        \operatorname{SynTO}_{\mathbb A,\mathbb B}(K)
    \]
    is the smallest regular quotient of
    \[
        \mathbb A^{\mathrm{op}}
    \]
    through which the residual Kleisli action on
    \[
        D=\operatorname{Der}_0(K)
    \]
    factors.
\end{corollary}

\begin{remark}[Quotient order]
\label{rem:syn-to-smallest-quotient-order}
    Here ``smallest'' is meant in the quotient order: every regular quotient of
    \(\mathbb A^{\mathrm{op}}\) through which \(\chi_D\) factors maps
    regular-epimorphically onto
    \[
        \operatorname{SynTO}_{\mathbb A,\mathbb B}(K).
    \]
\end{remark}

\begin{theorem}[Context expansion]
\label{thm:syn-to-context-expansion}
    Let
    \[
        a,a':u\to v
    \]
    be parallel arrows of \(\mathbb A\). Then
    \[
        a\equiv_K^{\mathrm{syn}}a'
    \]
    if and only if, for every specified behaviour \(\kappa\in K\) and every
    residual context
    \[
        r:v\to p_K(\kappa),
    \]
    the following two equations hold:
    \[
        \partial_{r\circ a}k(\kappa)
        =
        \partial_{r\circ a'}k(\kappa),
    \]
    and
    \[
        k(\kappa)(a,r)=k(\kappa)(a',r).
    \]
    Here
    \[
        k(\kappa)(a,r):
        k(\kappa)(r\circ a)\to k(\kappa)(r)
    \]
    denotes the arrow assigned by the residual behaviour
    \[
        k(\kappa):\mathbb A/p_K(\kappa)\to\mathbb B
    \]
    to the residual triangle
    \[
        u\xrightarrow{a}v\xrightarrow{r}p_K(\kappa).
    \]
\end{theorem}

\begin{proof}
    \leavevmode
    \begin{enumerate}
        \item \textbf{Reduce to equality of Kleisli actions on \(D_v\).}
        By Proposition~\ref{prop:syn-to-congruence-normal-form},
        \[
            a\equiv_K^{\mathrm{syn}}a'
        \]
        if and only if the two induced Kleisli maps
        \[
            D_v\rightsquigarrow D_u
        \]
        are equal.

        \item \textbf{Cover the fibre \(D_v\) by formal derivatives.}
        Pull back the regular epimorphism
        \[
            e_K:D_K\twoheadrightarrow D
        \]
        along the fibre \(D_v\hookrightarrow D\). Since regular epimorphisms are
        stable under pullback, this gives a regular epimorphic cover of \(D_v\).
        A generalized element of this cover is a pair
        \[
            (r,\kappa)
        \]
        where
        \[
            r:v\to p_K(\kappa).
        \]
        It represents the residual state
        \[
            \partial_r k(\kappa)\in D_v.
        \]

        The cover is displayed by
        \[
        \begin{tikzcd}[column sep=large, row sep=large]
            (D_K)_v
                \arrow[r, two heads]
                \arrow[d]
                \arrow[dr, phantom, "\lrcorner", very near start]
            &
            D_v
                \arrow[d]
            \\
            D_K
                \arrow[r, two heads, "e_K"']
            &
            D .
        \end{tikzcd}
        \]

        \item \textbf{Compute the transition part on the cover.}
        Acting by \(a\) on the representative \((r,\kappa)\) gives
        \[
            \partial_a(\partial_r k(\kappa))
            =
            \partial_{r\circ a}k(\kappa),
        \]
        while acting by \(a'\) gives
        \[
            \partial_{a'}(\partial_r k(\kappa))
            =
            \partial_{r\circ a'}k(\kappa).
        \]
        Thus equality of the transition parts is exactly
        \[
            \partial_{r\circ a}k(\kappa)
            =
            \partial_{r\circ a'}k(\kappa).
        \]

        \item \textbf{Compute the output part on the cover.}
        The restricted Mealy structure on \(D\) is the restriction of the
        canonical behaviour Mealy morphism. Hence
        \[
            \omega_D(a,\partial_r k(\kappa))
            =
            k(\kappa)(a,r),
        \]
        and similarly
        \[
            \omega_D(a',\partial_r k(\kappa))
            =
            k(\kappa)(a',r).
        \]
        Therefore equality of output components is exactly
        \[
            k(\kappa)(a,r)=k(\kappa)(a',r).
        \]
        The first equality of derivative states ensures that this equality of
        output arrows is well-typed: the two arrows have the same source and the
        same target.

        \item \textbf{Display the residual context.}
        The relevant residual triangle is
        \[
        \begin{tikzcd}[column sep=large, row sep=large]
            u
                \arrow[r, "a"]
                \arrow[rr, "r\circ a", bend left=16]
            &
            v
                \arrow[r, "r"]
            &
            p_K(\kappa).
        \end{tikzcd}
        \]
        Replacing \(a\) by \(a'\) gives the comparison triangle for the second
        equation.
    \end{enumerate}
\end{proof}

\subsection{Running specializations}
\label{subsec:syn-to-context-running-specializations}

In the one-object case, the contextual form says that
\[
    m\equiv_K^{\mathrm{syn}}n
\]
if and only if, for every specified behaviour \(\kappa\in K\) and every
residual context \(r\in M\),
\[
    \partial_{rm}k(\kappa)
    =
    \partial_{rn}k(\kappa),
\]
and the emitted outputs agree:
\[
    k(\kappa)(m,r)=k(\kappa)(n,r).
\]
Thus the congruence is a two-sided contextual congruence: the input is tested
after all generated left residual contexts and before all future residual
continuations contained in the equality of behaviours.

In the stateless case, for a functor \(F:\mathbb A\to\mathbb B\), parallel
arrows
\[
    a,a':u\to v
\]
are identified only if every generated residual context sees the same residual
functor and the same arrow value under \(F\):
\[
    F_1(a)=F_1(a')
\]
in the relevant contextual residual fibre.

\section{State quotients, syntactic quotients, and recognizers}
\label{sec:syn-to-state-vs-syntactic-recognizers}

The quotient constructed in this chapter is distinct from the state quotient
constructed in Chapter~\ref{ch:behaviour-categories-myhill-nerode}. The
state-level Nerode quotient is
\[
    D_K/{\equiv_K}
    \cong
    \operatorname{Der}_0(K).
\]
It identifies formal residual derivatives
\[
    (r,\kappa),
    \qquad
    (r',\kappa')
\]
precisely when they determine the same represented residual behaviour:
\[
    \partial_r k(\kappa)
    =
    \partial_{r'}k(\kappa').
\]
Its quotient is the residual state object
\[
    D=\operatorname{Der}_0(K).
\]

The syntactic transition--output quotient is instead
\[
    \mathbb A^{\mathrm{op}}/{\equiv_K^{\mathrm{syn}}}
    \cong
    \operatorname{SynTO}_{\mathbb A,\mathbb B}(K).
\]
It identifies input arrows precisely when they determine the same Kleisli
operation on all generated residual states.

\begin{proposition}[State quotient versus syntactic quotient]
\label{prop:syn-to-state-vs-syntactic}
    The state quotient
    \[
        D_K/{\equiv_K}
    \]
    classifies generated residual behaviours. The syntactic quotient
    \[
        \mathbb A^{\mathrm{op}}/{\equiv_K^{\mathrm{syn}}}
    \]
    classifies input arrows by their residual Kleisli action on those generated
    behaviours.
\end{proposition}

\begin{proof}
    \leavevmode
    \begin{enumerate}
        \item \textbf{Interpret the state quotient.}
        The state quotient is the quotient of the residual evaluation map
        \[
            d_K:D_K\to\mathbf{Beh}_0.
        \]
        Hence it identifies formal residual expressions exactly when they have
        the same residual behaviour.

        \item \textbf{Interpret the syntactic quotient.}
        The syntactic quotient is the quotient of the residual Kleisli action
        representation
        \[
            \chi_D:
            \mathbb A^{\mathrm{op}}
            \to
            \operatorname{KlEnd}_{\mathbb B}(D).
        \]
        Hence it identifies input arrows exactly when they induce the same
        transition--output operation on every generated residual state.

        \item \textbf{Display the two quotients.}
        The distinction is summarized by
        \[
        \begin{tikzcd}[column sep=huge, row sep=large]
            D_K
                \arrow[r, two heads]
                \arrow[dr, "d_K"']
            &
            D_K/{\equiv_K}
                \arrow[d, hook, "j_K"']
            \\
            &
            \mathbf{Beh}_0
        \end{tikzcd}
        \qquad
        \begin{tikzcd}[column sep=huge, row sep=large]
            \mathbb A^{\mathrm{op}}
                \arrow[r, two heads]
                \arrow[dr, "\chi_D"']
            &
            \mathbb A^{\mathrm{op}}/{\equiv_K^{\mathrm{syn}}}
                \arrow[d, hook, "m_{\mathrm{syn}}"']
            \\
            &
            \operatorname{KlEnd}_{\mathbb B}(D).
        \end{tikzcd}
        \]
    \end{enumerate}
\end{proof}

\subsection{Recognizer quotients}
\label{subsec:syn-to-recognizer-quotients}

The syntactic transition--output category may also be obtained as a quotient
from any realization of the same specification.

Let
\[
    i:K\to X
\]
be a realization of
\[
    k:K\to\mathbf{Beh}_0.
\]
Thus
\[
    X:\mathbb A\rightsquigarrow\mathbb B
\]
is a Mealy morphism and
\[
    \beta_X i=k.
\]

Let
\[
    G:=\operatorname{Reach}_0(X,K)\hookrightarrow X
\]
be the reachable part of this realization. By the semantic quotient theorem of
Chapter~\ref{ch:behaviour-categories-myhill-nerode}, the behaviour map induces
a regular epimorphic strict semantic homomorphism
\[
    h:G\twoheadrightarrow D,
    \qquad
    D=\operatorname{Der}_0(K).
\]

The reachable Mealy morphism \(G\) has its own residual Kleisli action
\[
    \chi_G:
    \mathbb A^{\mathrm{op}}
    \to
    \operatorname{KlEnd}_{\mathbb B}(G).
\]

\begin{definition}[Recognizer-generated syntactic category]
\label{def:syn-to-recognizer-generated-syntactic-category}
    The \textbf{recognizer-generated syntactic transition--output category}
    associated to \(i:K\to X\) is
    \[
        \operatorname{SynTO}_X(K)
        :=
        \operatorname{RegIm}(\chi_G).
    \]
\end{definition}

\begin{theorem}[Recognizer quotient theorem]
\label{thm:syn-to-recognizer-quotient}
    For every realization
    \[
        i:K\to X
    \]
    there is a canonical regular epimorphism of internal categories
    \[
        \operatorname{SynTO}_X(K)
        \twoheadrightarrow
        \operatorname{SynTO}_{\mathbb A,\mathbb B}(K).
    \]
    If \(X\) is reachable and behaviourally separated, this morphism is an
    isomorphism.
\end{theorem}

\begin{proof}
    \leavevmode
    \begin{enumerate}
        \item \textbf{Compare residual actions through the semantic quotient.}
        Suppose \(a,a':u\to v\) satisfy
        \[
            \chi_G(a)=\chi_G(a')
        \]
        as Kleisli maps
        \[
            G_v\rightsquigarrow G_u.
        \]
        Let \(x\in D_v\). Since
        \[
            h_v:G_v\twoheadrightarrow D_v
        \]
        is a regular epimorphism, equality of the corresponding transition and
        output components on \(D_v\) may be checked after pulling back along
        \(h_v\). For \(z\in G_v\), strict preservation of transitions gives
        \[
            h(\partial_a^G z)=\partial_a^D h(z),
            \qquad
            h(\partial_{a'}^G z)=\partial_{a'}^D h(z).
        \]
        Since \(\chi_G(a)=\chi_G(a')\), the two left-hand transition states in
        \(G\) are equal. Therefore the two induced transition states in \(D\)
        are equal after precomposition with \(h_v\), and hence are equal.

        Strict preservation of outputs gives
        \[
            \omega_G(a,z)=\omega_D(a,hz),
            \qquad
            \omega_G(a',z)=\omega_D(a',hz).
        \]
        Since \(\chi_G(a)=\chi_G(a')\), the output arrows on \(G\) are equal.
        Therefore
        \[
            \omega_D(a,hz)=\omega_D(a',hz)
        \]
        after precomposition with the regular epimorphism \(h_v\), and hence the
        output components of \(\chi_D(a)\) and \(\chi_D(a')\) are equal.
        Consequently
        \[
            \chi_D(a)=\chi_D(a'),
        \]
        so
        \[
            \ker(\chi_G)\subseteq\ker(\chi_D).
        \]

        \item \textbf{Descend the syntactic quotient.}
        Because
        \[
            \ker(\chi_G)\subseteq\ker(\chi_D),
        \]
        the image quotient for \(\chi_D\) coequalizes the kernel pair of the
        image quotient for \(\chi_G\). By exactness, it descends uniquely to a
        regular epimorphism
        \[
            \operatorname{SynTO}_X(K)
            \twoheadrightarrow
            \operatorname{SynTO}_{\mathbb A,\mathbb B}(K).
        \]

        This is displayed by
        \[
        \begin{tikzcd}[column sep=huge, row sep=large]
            \mathbb A^{\mathrm{op}}
                \arrow[r, two heads]
                \arrow[d, two heads]
            &
            \operatorname{SynTO}_X(K)
                \arrow[d, dashed, two heads]
            \\
            \operatorname{SynTO}_{\mathbb A,\mathbb B}(K)
                \arrow[r, equal]
            &
            \operatorname{SynTO}_{\mathbb A,\mathbb B}(K).
        \end{tikzcd}
        \]

        \item \textbf{Handle the reachable separated case.}
        If \(X\) is reachable, then
        \[
            G\cong X.
        \]
        If \(X\) is also behaviourally separated, the minimal realization theorem
        of Chapter~\ref{ch:behaviour-categories-myhill-nerode} gives an
        isomorphism of Mealy morphisms
        \[
            X\cong D.
        \]
        Hence the residual Kleisli actions \(\chi_G\) and \(\chi_D\) agree up
        to isomorphism, and their regular images are isomorphic.
    \end{enumerate}
\end{proof}

\subsection{Running specializations}
\label{subsec:syn-to-state-recognizer-running-specializations}

For an ordinary deterministic Mealy recognizer with reachable state set \(G\)
and minimal residual state set \(D\), the map
\[
    G\twoheadrightarrow D
\]
identifies behaviourally equivalent reachable states. The syntactic
transition--output monoid computed inside \(G\) maps onto the syntactic
transition--output monoid computed on \(D\). If \(G\) is already reachable and
behaviourally separated, these monoids coincide.

In the stateless case, a recognizer quotient compares the transition--output
category generated by a concrete functorial presentation with the intrinsic
syntactic transition--output category generated by its residual behaviours.

\section{Classical recovery and examples}
\label{sec:syn-to-classical-recovery-examples}

We now recover the ordinary syntactic monoid of a language as a special case.

\begin{theorem}[Classical syntactic monoid recovery]
\label{thm:syn-to-classical-recovery}
    Let
    \[
        L\subseteq\Sigma^\ast
    \]
    be a language. Let \(\mathbb A\) be the one-object category associated to
    the free monoid \(\Sigma^\ast\), and let \(\mathbb B\) be the chaotic
    category on
    \[
        2=\{0,1\}.
    \]
    Take \(K=1\), and let
    \[
        k:1\to\mathbf{Beh}_0
    \]
    classify the residual behaviour
    \[
        \lambda_L:\Sigma^\ast\to 2
    \]
    given by
    \[
        \lambda_L(w)=1
        \quad\Longleftrightarrow\quad
        w\in L.
    \]
    Then
    \[
        \operatorname{SynTO}_{\Sigma^\ast,\mathbb B}(L)
        \cong
        \operatorname{Syn}(L),
    \]
    the classical syntactic monoid of \(L\).
\end{theorem}

\begin{proof}
    \leavevmode
    \begin{enumerate}
        \item \textbf{Identify residual states.}
        A residual behaviour
        \[
            \Sigma^\ast/{*}\to\mathbb B
        \]
        is determined by its object assignment
        \[
            \lambda:\Sigma^\ast\to 2,
        \]
        because \(\mathbb B\) is chaotic. For \(u\in\Sigma^\ast\), the
        derivative is
        \[
            (\partial_u\lambda_L)(w)=\lambda_L(uw).
        \]
        Hence the derivative image is the ordinary derivative state set
        \[
            D(L)=\{\partial_uL\mid u\in\Sigma^\ast\}.
        \]

        \item \textbf{Use the context expansion theorem.}
        Since \(\mathbb B\) is chaotic, output comparisons are uniquely
        determined by their endpoints. Therefore equality of
        transition--output operations is equality of the induced state
        transformations on \(D(L)\).

        By Theorem~\ref{thm:syn-to-context-expansion}, two words
        \(u,v\in\Sigma^\ast\) are syntactically equivalent precisely when, for
        every residual context \(r\in\Sigma^\ast\),
        \[
            \partial_{ru}L=\partial_{rv}L.
        \]

        \item \textbf{Expand derivatives.}
        The equality
        \[
            \partial_{ru}L=\partial_{rv}L
        \]
        means that for every continuation \(s\in\Sigma^\ast\),
        \[
            rus\in L
            \quad\Longleftrightarrow\quad
            rvs\in L.
        \]
        Thus
        \[
            u\equiv v
            \quad\Longleftrightarrow\quad
            \forall r,s\in\Sigma^\ast,\;
            rus\in L\Longleftrightarrow rvs\in L.
        \]
        This is exactly the classical syntactic congruence of \(L\). Therefore
        the quotient is the classical syntactic monoid.

        \item \textbf{Display the two-sided context.}
        The classical context expansion is displayed by
        \[
        \begin{tikzcd}[column sep=large, row sep=large]
            *
                \arrow[r, "s"]
                \arrow[rr, "us", bend left=16]
                \arrow[rrr, "rus", bend left=30]
            &
            *
                \arrow[r, "u"]
            &
            *
                \arrow[r, "r"]
            &
            * .
        \end{tikzcd}
        \]
        The equality condition compares this diagram with the corresponding
        diagram obtained by replacing \(u\) with \(v\).
    \end{enumerate}
\end{proof}

\begin{example}[Deterministic Mealy machines]
\label{ex:syn-to-deterministic-mealy-machines}
    Let
    \[
        M=I^\ast
    \]
    be a free input monoid and let \(N\) be an output monoid. A deterministic
    Mealy behaviour has a derivative state object \(D\) and a residual cocycle
    \[
        \omega_D:I^\ast\times D\to N.
    \]
    The syntactic transition--output monoid is
    \[
        \operatorname{im}
        \left(
            I^\ast\to D^D\ltimes N^D
        \right).
    \]
    A word \(w\in I^\ast\) acts by
    \[
        x\longmapsto x\cdot w
    \]
    and emits
    \[
        x\longmapsto \omega_D(w,x).
    \]
    Thus the syntactic monoid records both the residual state transformation and
    the cumulative output emitted from each residual state.
\end{example}

\begin{example}[Group-valued cumulative outputs]
\label{ex:syn-to-group-valued-cumulative-outputs}
    Let \(G\) be a group and let
    \[
        \ell:\Sigma^\ast\to G
    \]
    be a cumulative output function. It determines an incremental residual
    cocycle by
    \[
        \lambda_\ell(r,m)=\ell(r)^{-1}\ell(rm).
    \]
    The syntactic transition--output congruence identifies two input words
    \(m,n\) precisely when every residual context \(r\) sees the same future
    residual behaviour and the same group increment:
    \[
        \partial_{rm}\lambda_\ell=\partial_{rn}\lambda_\ell,
        \qquad
        \lambda_\ell(r,m)=\lambda_\ell(r,n).
    \]
\end{example}

\begin{example}[Base-\texorpdfstring{\(r\)}{r} carry machine]
\label{ex:syn-to-base-r-carry-machine}
    Let
    \[
        M=N=(\mathbb N,+),
        \qquad
        P=\{0,\dots,r-1\}.
    \]
    The carry machine has transition
    \[
        x\cdot n=(x+n)\bmod r
    \]
    and output
    \[
        \omega(n,x)=
        \left\lfloor\frac{x+n}{r}\right\rfloor.
    \]
    Its syntactic transition--output monoid is the image of
    \[
        \mathbb N\to P^P\ltimes\mathbb N^P.
    \]
    The element \(n\in\mathbb N\) acts by the residue transformation
    \[
        x\longmapsto (x+n)\bmod r
    \]
    and emits the carry function
    \[
        x\longmapsto
        \left\lfloor\frac{x+n}{r}\right\rfloor.
    \]
\end{example}

\begin{example}[Typed path-category machines]
\label{ex:syn-to-typed-path-category-machines}
    Let \(G\) and \(H\) be directed graphs, and consider a Mealy morphism
    \[
        \mathbf{Path}(G)\rightsquigarrow\mathbf{Path}(H).
    \]
    The residual state object is indexed by vertices of \(G\). A path
    \[
        a:u\to v
    \]
    induces a Kleisli map
    \[
        D_v\rightsquigarrow D_u
    \]
    whose state part is residual reindexing along \(a\), and whose output part
    is the emitted path in \(\mathbf{Path}(H)\). The syntactic
    transition--output category is the image of
    \[
        \mathbf{Path}(G)^{\mathrm{op}}
    \]
    in the corresponding Kleisli endomorphism category.
\end{example}

\begin{example}[Register allocation]
\label{ex:syn-to-register-allocation}
    In the register-allocation Mealy morphism, a state is a bijective
    allocation
    \[
        \rho:\Gamma\cong R.
    \]
    For an inclusion
    \[
        \Gamma'\hookrightarrow\Gamma,
    \]
    the residual transition restricts the allocation to
    \[
        \rho|_{\Gamma'}:\Gamma'\cong \rho(\Gamma'),
    \]
    and the emitted output is the inclusion
    \[
        \rho(\Gamma')\hookrightarrow R.
    \]
    The syntactic transition--output category records the generated restriction
    operations on residual allocation states together with the emitted
    register-interface inclusions.

    The transition--output square is
    \[
    \begin{tikzcd}[column sep=large, row sep=large]
        \Gamma'
            \arrow[r, hook]
            \arrow[d, "{\rho|_{\Gamma'}}"', "\sim"]
        &
        \Gamma
            \arrow[d, "{\rho}", "\sim"']
    \\
        \rho(\Gamma')
            \arrow[r, hook]
        &
        R .
    \end{tikzcd}
    \]
\end{example}

\begin{example}[Stateless functors]
\label{ex:syn-to-stateless-functors}
    Let
    \[
        F:\mathbb A\to\mathbb B
    \]
    be an internal functor, regarded as a map-representable Mealy morphism. Its
    residual behaviours are
    \[
        H_v^F=F\circ\operatorname{dom}_v.
    \]
    For an arrow
    \[
        a:u\to v,
    \]
    the residual Kleisli action sends
    \[
        H_v^F\longmapsto H_u^F
    \]
    and emits the arrow
    \[
        F_1(a):F_0(u)\to F_0(v).
    \]
    Thus the syntactic transition--output category is the image of
    \[
        \mathbb A^{\mathrm{op}}
    \]
    acting on the residual restrictions of \(F\), with output arrows given by
    the arrow map of \(F\).
\end{example}

\section{Chapter summary}
\label{sec:syn-to-summary}

Starting from a behaviour specification
\[
    k:K\to\mathbf{Beh}_0,
\]
Chapter~\ref{ch:behaviour-categories-myhill-nerode} constructs the derivative
image
\[
    D=\operatorname{Der}_0(K)
\]
and the minimal residual Mealy morphism
\[
    \mathsf{Min}(K):\mathbb A\rightsquigarrow\mathbb B.
\]
The present chapter studies the input operations induced on \(D\).

The output category \(\mathbb B\) determines the output-comparison monad
\[
    T_{\mathbb B}=(t_B)_!(s_B)^*
\]
on
\[
    \mathcal E/B_0.
\]
Its Kleisli category records state maps equipped with output-comparison arrows.
The residual carrier
\[
    D\to A_0\times B_0
\]
is an \(A_0\)-indexed object of this Kleisli category.

The Mealy structure on \(D\) gives a residual Kleisli action
\[
    \chi_D:
    \mathbb A^{\mathrm{op}}
    \longrightarrow
    \operatorname{KlEnd}_{\mathbb B}(D).
\]
The syntactic transition--output category is the regular image of this action:
\[
    \operatorname{SynTO}_{\mathbb A,\mathbb B}(K)
    =
    \operatorname{RegIm}(\chi_D).
\]
It is the smallest regular quotient of
\[
    \mathbb A^{\mathrm{op}}
\]
through which the residual Kleisli action on \(D\) factors.

The kernel pair of \(\chi_D\) is the syntactic transition--output congruence
\[
    {\equiv_K^{\mathrm{syn}}}.
\]
Exactness gives the quotient presentation
\[
    \mathbb A^{\mathrm{op}}/{\equiv_K^{\mathrm{syn}}}
    \cong
    \operatorname{SynTO}_{\mathbb A,\mathbb B}(K).
\]

The context expansion theorem identifies this congruence explicitly: two
parallel input arrows are syntactically equivalent precisely when every
specified residual context sees the same generated derivative and the same
emitted output. In the ordinary Boolean free-input case, where the output
category is chaotic, this recovers the classical syntactic monoid of a
language.

The constructions may be summarized as follows:
\[
\begin{array}{c|l}
\text{structure} & \text{role} \\
\hline
T_{\mathbb B}=(t_B)_!(s_B)^*
&
\text{output-comparison monad}
\\
\operatorname{Kl}_{\mathbb B}
&
\text{Kleisli category of output comparisons}
\\
D=\operatorname{Der}_0(K)
&
\text{minimal residual state object}
\\
\operatorname{KlEnd}_{\mathbb B}(D)
&
\text{Kleisli endomorphism category of the residual family}
\\
\chi_D:\mathbb A^{\mathrm{op}}\to\operatorname{KlEnd}_{\mathbb B}(D)
&
\text{residual transition--output action}
\\
\operatorname{SynTO}_{\mathbb A,\mathbb B}(K)
&
\text{regular image of the residual action}
\\
{\equiv_K^{\mathrm{syn}}}
&
\text{kernel-pair syntactic transition--output congruence}
\\
\mathbb A^{\mathrm{op}}/{\equiv_K^{\mathrm{syn}}}
&
\text{syntactic transition--output quotient}
\\
\operatorname{SynTO}_X(K)\twoheadrightarrow
\operatorname{SynTO}_{\mathbb A,\mathbb B}(K)
&
\text{recognizer-generated quotient}
\end{array}
\]